\section{Performance Measuring Metrics}\label{sec:metrics}

\paragraph{Throughput (GFLOP/s).} SpMV performs $2\,\mathrm{nnz}$ flops per call and SpGEMM two flops per
multiply-add. GFLOP/s normalises time by the work done, so matrices of different sizes can be compared, and it is the
unit both papers report, which lets us compare with them directly.

\paragraph{Effective bandwidth (fraction of measured peak).} SpMV and SpGEMM are memory-bound, so the fraction of the
\emph{measured} \stat{pf-device-read}\,GB/s a kernel reaches says how close it is to the hardware limit, and therefore
how much any other format could still gain. We use the measured rather than the datasheet peak (192\,GB/s), which no
kernel can reach. Nsight Compute's DRAM, L2 and occupancy counters explain \emph{why} a kernel sits where it does.

\paragraph{Slowdown against the oracle, $T/T_\text{oracle}$.} The oracle is the fastest of all configurations for
each matrix, found by exhaustive measurement. The ratio says how much of the attainable speed a selector leaves
behind, independently of how good cuSPARSE happens to be on that matrix. We aggregate ratios with the \emph{geometric}
mean, the correct mean for ratios (a 2\x{} gain and a 2\x{} loss cancel), and also report the worst case, since a
selector that is good on average but occasionally very slow is not safe to deploy.

\paragraph{Performance profiles.} Dolan--Mor\'e profiles~\cite{dolan2002profiles} plot, for every factor $\tau$, the
fraction of matrices on which a method is within $\tau$ of the oracle. They show the whole distribution, not only its
mean, and are the standard way to compare solvers over a test set.

\paragraph{Selection overhead in SpMV equivalents.} Feature extraction and conversion times are divided by the time of
one cuSPARSE SpMV on the same matrix. This expresses an overhead in the unit that matters to the caller: how many calls
it takes to pay it back.

\paragraph{Total time for $N$ calls.} Plan cost plus $N$ SpMVs, relative to always using cuSPARSE, for $N$ from 1 to
1000. This is the quantity break-even optimises, and it makes the one-call case ($N=1$) explicit.

\paragraph{Speedup.} Against Algorithm~4 for out-of-core SpGEMM (the paper's own baseline), against cuSPARSE for the
selector and CG, and against the sequential CPU code on the same machine for the GPU and OpenMP versions. For
OpenMP we also report speedup over the same code on one thread and parallel efficiency, which separate parallel
scaling from compiler or algorithm effects.

\paragraph{Data transfer cost.} PCIe bandwidth (host to device, device to host, both), the slowdown when an operand
lives in host memory, and out-of-core throughput as a fraction of the device-only kernel's. These quantify how much the
host link costs, which is the central constraint of out-of-core work.

\paragraph{Time to solution.} For CG and the AMG set-up we time the whole application, including planning and
transfers, because a faster kernel only matters if it shortens the end result.

\paragraph{Cost-model accuracy.} The ratio of the planner's predicted execution time to the measured one shows whether
the simulated pipeline is trustworthy enough to choose plans with.
